\subsection{Authentication and Identity}
\label{sec:impl_auth}

The identity subsystem handles user registration, Google OAuth sign-in, JWT session management, and multi-factor authentication (MFA). It is implemented in \texttt{abridgeai/features/identity/} and exposes four router groups: \texttt{/auth}, \texttt{/me}, \texttt{/users} (admin), and \texttt{/mfa}.

\paragraph{Google OAuth flow}
The platform uses Google OAuth 2.0 as its sole external identity provider. The sign-in flow is implemented in \texttt{features/identity/services/login.py} and proceeds as follows:

\begin{enumerate}
    \item The frontend redirects the user to Google's authorisation endpoint with the configured \texttt{client\_id} and \texttt{redirect\_uri}.
    \item Google redirects back with an authorisation code, which the frontend forwards to \texttt{POST /auth/google/callback}.
    \item The backend exchanges the code for a Google profile via \texttt{infrastructure/google\_oauth.py}, which calls Google's \texttt{tokeninfo} endpoint.
    \item A pre-registration gate checks that the email already exists in the \texttt{users} table with \texttt{status = 'active'}. New OAuth profiles from unregistered emails are rejected with HTTP 403, keeping the platform invite-only.
    \item On success, an \texttt{AuthIdentity} row linking the user to the Google provider subject is created (or reused if it already exists), and a new \texttt{AuthSession} row is written with a hashed refresh token.
    \item The response returns a short-lived JWT access token (configurable TTL, default 15 minutes) and a long-lived refresh token.
\end{enumerate}

\paragraph{JWT and session management}
Access tokens are signed with HS256 using a secret loaded from the \texttt{JWT\_SECRET} environment variable. The \texttt{core/security} module provides \texttt{create\_access\_token()}, \texttt{hash\_secret()}, and \texttt{generate\_token()} (32-byte URL-safe random). Refresh tokens are stored in \texttt{auth\_sessions} as keyed HMAC-SHA256 digests --- the plaintext is never persisted --- and compared in constant time. A keyed MAC rather than a password hash is the right primitive here: the input is a 32-byte random token with no guessable structure, so the work factor that protects a human-chosen password buys nothing, while the server-held key means a leaked database alone does not yield usable refresh tokens. Token rotation is enforced: each \texttt{POST /auth/refresh} call invalidates the old session row and issues a new one.

\paragraph{Multi-factor authentication}
TOTP-based MFA is implemented in \texttt{features/identity/services/mfa.py}. Enrollment generates a TOTP secret, returns a provisioning URI for authenticator apps, and stores the encrypted secret only after the user verifies a valid code. Subsequent logins that have MFA enabled require a second \texttt{POST /mfa/verify} call before the access token is issued. The MFA gate is enforced at the router level via a FastAPI dependency that checks the \texttt{mfa\_verified} claim in the JWT.

\paragraph{Role-based access}
Every authenticated request resolves the caller's roles and permissions through the access control subsystem (Section~\ref{sec:impl_access_control}). The identity layer attaches the resolved \texttt{UserContext} — containing user ID, organisation memberships, and permission set — to the request state so downstream route handlers can perform permission checks without additional database queries.

\paragraph{The session lifecycle}
There is one way in and several ways out. Google OAuth is the only sign-in route --- the platform holds no passwords and has no local credential to verify --- so every session begins with the pre-registration gate above, and an unrecognised Google profile is refused before any session row exists. Where the account has MFA enrolled, the callback yields a token that carries no \texttt{mfa\_verified} claim and therefore opens no protected route; only the subsequent verification exchanges it for a usable access token. The second factor is thus a gate on issue, not a check applied afterwards.

A live session then persists until something ends it. Ordinary expiry is the common case, the access token being short-lived and each refresh rotating the session row so a replayed refresh token finds its row already invalidated. Explicit logout terminates the presented session. Disabling an account terminates all of them at once, which is why the reverse index from user to session IDs exists: revocation must not depend on scanning the session keyspace, and an administrator disabling an account expects it to take effect on the next request rather than at the next token expiry.
